\section{Modeling and Inference}
\label{sec:modeling}

% Rubric (Methodology): "Method fully specified and justified. Enough detail
% for replication."
% Guidelines:
%   - Specify the probabilistic model.
%   - Assumptions, parameters, dependency structure; relate to the data/problem.
%   - Inference / learning algorithms implemented.
%
% Outline:
\subsection{Bayesian Network}
\label{sec:bn}
% - Structure figure (nodes, edges); justify each edge.
% - Conditional independence assumptions and what they mean physically.
% - Parameters: CPTs; count of free parameters.
\todo{BN structure, assumptions, parameterization.}

\subsection{Learning with Missing Data (EM)}
\label{sec:em}
% - E-step / M-step for this model \citep{dempster1977em}; initialization,
%   restarts, convergence criterion.
\todo{EM derivation / description for our model.}

\subsection{Inference}
\label{sec:inference}
% - Queries we answer, e.g. P(true level | sensors, weather).
% - Algorithm (exact: variable elimination); complexity.
\todo{Inference queries and algorithm.}

\subsection{Temporal Extension: HMM}
\label{sec:hmm}
% - Hidden state, emissions, transition assumptions \citep{rabiner1989hmm}.
% - Baum--Welch for learning; forward--backward / Viterbi for inference.
\todo{HMM model and algorithms (if pursued).}
